\chapter{Discussion and Conclusion}\label{ch:conclusion}

\section{Core finding}

The final one-lead system reached high overall accuracy on later MLII ECG from
patients whose earlier MLII ECG was used for training. Whole-system pooled
accuracy was 96.86\%, equal-subject accuracy was 96.90\%, and the
subject-clustered 95\% confidence interval was 94.02--99.05\%. Its lower bound
is above 90\%, and two non-parametric subject-level tests also reject a 90\%
null hypothesis.

This is a genuine chronological holdout, not a shuffled-beat calculation. The
final 20\% was excluded from fitting and selection, and a 16-beat gap reduced
contact between neighbouring windows. However, it is a patient-adapted result:
the same patients appear in the first 80\% training region. It must not be
described as 96.86\% accuracy on unseen patients.

A limited post-freeze INCART check was added after the main final experiment.
It used lead II from an unseen source database and is useful for demonstrating
external-domain behaviour, including example ECG segments. It does not replace
a large patient-level external validation study.

\section{Were the goals reached?}

\begin{itemize}
    \item \textbf{Aggregate accuracy above 90\%: reached.} Abnormality,
    conditional subtype, whole-system accuracy, and whole-system macro F1 all
    exceed 90\%.
    \item \textbf{Whole-system 95\% interval above 90\%: reached.} The lower
    bound is 94.02\%.
    \item \textbf{Interval width at most three points: not reached.} The width
    is 5.03 points.
    \item \textbf{Every arrhythmia above 90\% recall: not reached.} PAC recall
    is 70.05\%.
    \item \textbf{More than 90\% accuracy on completely unseen patients:
    partially checked, not fully validated.} The limited INCART check exceeded
    90\% on its saved subset, but it is too small to establish patient-level
    external generalisation.
\end{itemize}

\section{Why the final architecture worked}

The system works as a one-channel system. The autoencoder uses only the first
stored channel of NSRDB, while every supervised and deployed input is MLII. The
NSRDB header calls its channel ECG1 but does not document an anatomical Lead-II
placement. For that reason, this thesis does not claim that the two source
channels have the same electrical axis. Removing MITDB records without MLII
also prevents the supervised system from silently changing viewpoints.

The model combines evidence that explains different parts of the ECG. Beat bins
and derivatives describe shape. Spectral features describe frequency content.
Causal RR history describes timing. Normal-only reconstruction describes
mismatch from healthy structure. Extra Trees learn non-linear links between
these features without assuming a straight-line decision boundary. Earlier
labelled ECG also gives a patient-specific reference, which makes persistent
bundle-branch and rhythm patterns easier to recognise later in the same
recording.

\section{Remaining weaknesses}

PAC remains difficult because it can look like normal rhythm or PVC and often
depends on subtle timing evidence. The final system correctly classified 414
of 591 PAC beats end to end. Its PAC precision is high, but sensitivity is low.
Improvement requires more independent PAC patients and a larger external test,
not more tuning against the temporal test or the limited INCART check already
reported here.

The interval width is controlled by 45 independent subject clusters, including
three below-target clusters. Adding more beats from the same records would
produce a misleadingly narrow beat-level interval without adding independent
evidence. A genuinely shorter interval requires more independent patients with
enough examples from all classes.

\section{Limitations}

\begin{enumerate}
    \item The temporal study has patient overlap by design and cannot establish
    cold-start patient generalisation.
    \item The supervised data represent 45 MITDB subject clusters; no external
    hospital cohort was used. The added INCART check is an external database
    check, not a hospital-cohort validation study.
    \item The healthy autoencoder used complete recordings from all 18 NSRDB
    subjects. Its first stored channel is named ECG1 in the WFDB headers, but
    the electrode placement is not documented as Lead II.
    \item The independent benefit of the four autoencoder features was not
    tested with a final held-out ablation.
    \item PAC end-to-end recall is 70.05\%.
    \item Unsupported abnormal symbols remain in the binary result but are
    excluded from the supported six-class result.
    \item Records 102 and 104 were excluded because they contain V5 and V2,
    not MLII. Record 114 uses its stored MLII channel by signal name rather
    than assuming a fixed channel number.
    \item Offline evaluation uses annotated R peaks. Live R-peak detection can
    add errors that are not included in the offline temporal result.
    \item Evaluation is beat based, not a clinical episode or patient-level
    diagnosis study.
\end{enumerate}

These limits follow the principle that development performance and clinical
evidence must be reported separately \cite{collins2024tripodai}.

\section{Required next study}

A future study should add many independent patients, reserve an external
patient-level database for one final test, pre-register a lower 95\% bound
above 90\%, keep the three-point width target, and set minimum recall for each
arrhythmia. All thresholds must be selected without the external test. The
study should also compare the full model with and without autoencoder
residuals, and validate real-time R-peak detection and Lead-II hardware
placement.

\section{Final statement}

The thesis demonstrates a one-lead, patient-adapted temporal ECG system with
96.86\% whole-system accuracy and a 94.02--99.05\% subject-clustered 95\%
interval. This is meaningful evidence for later MLII ECG from known patients.
The added INCART result is a useful external-source check, but it is not
evidence of equivalent performance for all new patients, and it does not solve
PAC sensitivity. The aggregate final-study goal is reached with clearly stated
statistical and clinical limits. The system remains for research use only.
